% Generated draft: requires mathematical and source review.
\documentclass[11pt]{article}
\usepackage[margin=25mm]{geometry}
\usepackage{fontspec}
\setmainfont{Latin Modern Roman}
\usepackage{amsmath,amssymb,mathtools,unicode-math}
\setmathfont{Latin Modern Math}
\usepackage{enumitem,needspace,xurl,hyperref}
\setlength{\parindent}{0pt}
\setlength{\parskip}{5pt}
\setlength{\emergencystretch}{3em}
\newcommand{\N}{\mathbb{N}}
\newcommand{\Z}{\mathbb{Z}}
\newcommand{\Q}{\mathbb{Q}}
\newcommand{\R}{\mathbb{R}}
\newcommand{\C}{\mathbb{C}}
\newcommand{\F}{\mathbb{F}}
\newcommand{\A}{\mathbb{A}}
\newcommand{\PP}{\mathbb{P}}
\newcommand{\E}{\mathbb{E}}
\newcommand{\eps}{\varepsilon}
\newcommand{\dd}{\,\mathrm{d}}
\begin{document}
\section*{UnsolvedMath ID: 701332}
\textbf{Source problem code:} catalogue entry
\subsection{Definitions and mathematical statement}
A group \(G\) is finitely presented if it admits a presentation with finitely many generators and finitely many defining relations. It is soluble if its derived series, defined by \(G^{(0)}=G\) and \(G^{(i+1)}=[G^{(i)},G^{(i)}]\), reaches the trivial group after finitely many steps. A group \(K\) is nilpotent if its lower central series \(\gamma_1(K)=K\), \(\gamma_{i+1}(K)=[\gamma_i(K),K]\), reaches the trivial group after finitely many steps.

Use the normal-series convention for nilpotent length: for a soluble group \(H\), let \(\ell_{\mathrm{nil}}(H)\) be the least integer \(n\geq 0\) for which there is a series
\[
1=H_0\leq H_1\leq\cdots\leq H_n=H,
\qquad H_i\trianglelefteq H,
\]
with every factor \(H_i/H_{i-1}\) nilpotent. In particular, \(\ell_{\mathrm{nil}}(1)=0\). This length is finite for soluble groups, since a finite derived series has abelian factors. The supplied source does not give a separate definition of nilpotent length; the convention above should be checked against the original source during review.

\textbf{Question (Bludov--Groves, Kourovka Notebook 17.30).} Does there exist an integer \(l\geq 0\) such that for every finitely presented soluble group \(G\) there exists a subgroup \(H\leq G\) satisfying
\[
[G:H]<\infty
\qquad\text{and}\qquad
\ell_{\mathrm{nil}}(H)\leq l?
\]
The constant \(l\) must be independent of \(G\). The subgroup \(H\) is not required to be normal, and no uniform bound on its index or on the nilpotency classes of the factors is requested. The source adds the comparison \emph{``Cf. 16.35 in Archive.''} The referenced archive statement is not supplied, so its relationship to this question remains unspecified.
\subsection{Short English statement}
Is there a universal bound on the nilpotent length of a suitable finite-index subgroup of every finitely presented soluble group? The supplied registry labels this question open; that label is not an independently verified assessment.
\subsection{Sources}
\begin{itemize}
\item Local UnsolvedMath registry, supplied record with permanent ID \(701332\), exact source problem code \texttt{catalogue entry}, title \emph{Kourovka Notebook Problem 17.30}. The supplied status is \texttt{open}.
\item Attributed to V. V. Bludov and J. R. J. Groves, proposed in 2010. Supplied source: \emph{The Kourovka Notebook (Unsolved Problems in Group Theory)}, problem 17.30, 17th issue (2010); the record also identifies the 21st issue edition (2026), \texttt{arXiv:1401.0300v48}. The supplied collection description names E. I. Khukhro and V. D. Mazurov as editors and covers Issues 1--21 (1965--2026). These bibliographic details are supplied attribution and have not been independently verified.
\item Supplied cross-reference: \emph{``Cf. 16.35 in Archive.''} No text or separate bibliographic details for that archive entry are supplied.
\end{itemize}
\end{document}
